\documentstyle[%


righttag%


,11pt%


,amssymb%


,russian%               remove in Eng.


,izvru%                 Set izven% in Eng.


]{amsart}





% {\ne}


\newcommand{\up}[2]{\overset{#1}{#2}{}}


\newcommand{\Lo}{\Longrightarrow}


\newcommand{\lo}{\longrightarrow}


\newcommand{\stac}{\triangleq}


\newcommand{\fo}{\forall}


\newcommand{\exi}{\exists}


\newcommand{\emp}{\emptyset}


\newcommand{\Gm}{\Gamma}


\newcommand{\pr}{\operatorname{pr}}


\newcommand{\Def}{\operatorname{def}}


\newcommand{\na}{\operatorname{na}}


\newcommand{\ass}{\operatorname{ass}}


\newcommand{\isot}{\operatorname{isot}}


\newcommand{\Isot}{\operatorname{Isot}}


\newcommand{\DOM}{\operatorname{DOM}}


\newcommand{\vt}{\vartheta}


\newcommand{\Om}{\Omega}


\newcommand{\vp}{\varphi}


\numberwithin{equation}{section}


\newcommand{\tts}[1]{}





\theoremstyle{plain}


\theorembodyfont{\it}


\newtheorem{thms}{\hskip\parindent Теорема}[section]


\newtheorem{props}{\hskip\parindent Предложение}[section]





\theoremstyle{definition}


\theorembodyfont{\rm}


\newtheorem{cors}{\hskip\parindent Следствие}[section]


\newtheorem{notes}{\hskip\parindent Замечание}[section]


\newtheorem{conds}{\hskip\parindent Условие}[section]





%\hoffset=-15mm%                  For Eng.


\setcounter{page}{77}


\god{2001}


\nomer{12~(475)} %                        Remove in Eng.


%\nomer{Vol.\,45, No.\,12}%               Set in Eng.


%\str{\pageref{begin}--\pageref{end}}%  Set in Eng.


\udk{517.972.8}





\author{\it А.Г.\,Ченцов}


% NB:   ^^ remove in Eng.


\title{К вопросу о двойственности различных версий\\ метода


программных итераций}


\thanks{Работа выполнена при финансовой поддержке Российского фонда


фундаментальных исследований (гранты \No\,01-01-96450, \No\,0001 00348)


и Международного научно-технического центра (проект \No\,1293).}





\date{}


%\pagestyle{myheadings}%         Set in Eng.


\pagestyle{plain} %              Remove in Eng.


\begin{document}


%\thispagestyle{plain}%          Set in Eng.


\maketitle


\label{begin}








Для аксиоматически определяемой наследственной системы


рассматриваются две конструкции, восходящие к


известному в теории дифференциальных игр методу программных


итераций: прямая (в смысле построения управляющих процедур) версия


и непрямая версия,


доставляющая объект, важный


для соответствующей задачи управления (функция цены, стабильный мост),


но являющийся посредником при построении управляющей процедуры.


Установлена двойственность этих


методов, охватывающая операторы, формирующие итерационные


последовательности, сами эти последовательности, их пределы и условия


невырожденности


обеих итерационных процедур.





\section{Введение}





Отметим некоторые сокращения, используемые ниже:


ДИ (дифференциальная игра), МПИ (метод программных итераций),


МО (многозначное отображение), м/с (мультиселектор), НМ


(направленное множество),


ОПП (оператор


программного поглощения), п/м (подмножество), п/п (подпространство),


ТП (топологическое пространство).





В теории ДИ используются методы,


связанные с построением стабильных мостов в смысле


Н.Н.\,Красовского и функции цены соответствующей ДИ


[1]--[5].


Известные трудности решения ДИ на основе динамического


программирования отмечены в [6]. Развитие этого метода


для решения ДИ потребовало применения негладкого


анализа в связи с негладкостью


функции цены [4]. В т.\,н. регулярных ДИ


применялись вспомогательные программные конструкции,


использующие игровые аналоги принципа максимума


Л.С.\,Понтрягина [1], [2].


Для общего случая нелинейной ДИ


Н.Н.\,Красовским и А.И.\,Субботиным была установлена фундаментальная теорема


об альтернативе [2], имеющая много важных следствий


[2]--[4]. Естественным образом возник


метод


программных итераций (МПИ) [7]--[11] решения ДИ в общем случае.


Обзор основных конструкций МПИ приведен в [5] и в [12].


Аналоги МПИ использовались в работах по теории


ДИ (отметим только некоторые наиболее близкие: [13]--[17]).


Кроме того, аналоги МПИ стали применяться для исследования


динамических задач другой природы: построение обобщенных решений


уравнения Гамильтона--Якоби [18], [19], построение неупреждающих


или наследственных м/с


МО в [20]--[23]. В конструкциях МПИ все более существенной


становилась методология, связанная с теоремами о неподвижной точке.


Отметим, в частности, работы Л.В.\,Канторовича в связи с конструкциями,


подобными ([24], c.\,237--239). Возникает вопрос о сравнении различных


версий МПИ: варианта [8],


реализуемого на пространстве п/м пространства


позиций, и варианта [20]--[23], именуемого ниже прямым.





В разделе 2 дана сводка обозначений


общего характера. Третий раздел содержит


более специальные понятия, связанные с представлением


абстрактной динамической системы, и обсуждение


некоторых возможных конкретизаций. Существенную роль


играет серия условий, используемых в разных сочетаниях


при построении непрямой версии МПИ, которая


излагается в разделе 4 (речь идет о существенном обобщении


результатов [8], [9]; последние определяют итерационный


метод, не доставляющий управляющую процедуру


непосредственно); главный результат раздела --- теорема 4.1 ---


определяет достаточное условие сходимости процедуры к


экстремальной неподвижной точке ОПП. Последний на идейном уровне


согласуется с аналогичным оператором [9] и отвечает ситуации,


когда произвольно выбираемая из заданного множества


помеха парируется полезным воздействием, которое формируется


в ответ на эту помеху и ориентируется на достижение нужного


качества (см. в этой связи обсуждение в [5], гл.\,IV). В пятом разделе


введен вариант прямой версии МПИ [20]--[23] и устанавливается


связь этой прямой и ранее излагаемой непрямой версий, имеющая смысл


двойственности: одна итерационная процедура полностью определяется


другой (см. теорему 5.1 и предшествующие положения).


В шестом разделе эта двойственность распространяется на


условия невырожденности, важные с точки зрения итерационного


построения неупреждающих откликов


на помеховые


реализации


(квазистратегий), используемых в качестве процедур управления.


Само появление итерационного процесса в разделе 4 связано,


как и в [8], [9], с игровым характером решаемой задачи, что проявляется


уже в определении ОПП и подчеркивается введением представлений на основе


поточечно исполняемого пересечения идемпотентных (в естественных


случаях) операторов.





\section{Общие определения и обозначения}





Используем символы


$\Def$ (по определению) и $\stac$ (равно по определению).


Через ${{\cal P}(H)}$ (${{\cal P}'(H)}$) обозначаем семейство


всех (всех непустых) п/м множества $H$. Через $B^A$ обозначаем


множество всех операторов из множества $A$ в множество $B$;


при $f \in {B^A}$ и $C \in {{\cal P}(A)}$ полагаем,


что $(f \mid C)$ есть сужение $f$ на $C$ ([25], c.\,26).


Множество, все элементы которого сами являются множествами,


называем семейством. Принимаем аксиому выбора. В дальнейшем


${\Bbb R}$ --- вещественная прямая, ${\cal N} \stac \{1;2;\dotsc\}$,


${{\cal N}_0} \stac \{0\} \cup {\cal N}$.


Через ${\bold I}_U$ обозначаем тождественное отображение


множества $U$ на себя: ${{\bold I}_U} \in {U^U}$ и при этом


$\fo{u} \in U:{{\bold I}_U}(u) \stac u$; если


$\alpha \in {U^U}$, то $({{\alpha}^k})_{k \in {{\cal N}_0}}$ есть


последовательность в $U^U$, для которой ${{\alpha}^0} \stac


{{\bold I}_U}$ и $\fo{s} \in {\cal N}:{{\alpha}^s} =


\alpha \circ {{\alpha}^{s-1}}$. Если ${\cal S}$ --- семейство всех п/м


некоторого множества и $\alpha \in {\cal S}^{\cal S}$, то


$\up{\infty}{\alpha} \in {\cal S}^{\cal S}$


действует по правилу: при $H \in {\cal S}$


множество $\up{\infty}{\alpha}(H)$ есть $\Def$ пересечение


всех множеств ${{\alpha}^k}(H)$, $k \in {{\cal N}_0}$. Если


$A$ и $B$ --- множества, то ${\Bbb M}(A,B) \stac {{\cal P}(B)}^A$;


для $\beta \in {{\Bbb M}(A,B)}^{{\Bbb M}(A,B)}$ определяем


${{\beta}^\infty} \in {{\Bbb M}(A,B)}^{{\Bbb M}(A,B)}$ по правилу:


при ${\cal C} \in {\Bbb M}(A,B)$ и $a \in A$ множество


${{\beta}^\infty}({\cal


C})(a)$ есть $\Def$ пересечение всех множеств


${{\beta}^k}({\cal C})(a)$, $k \in {{\cal N}_0}$. Далее используются


традиционные варианты монотонной сходимости.


Если $E$ --- множество, $(A_i)_{i \in {\cal N}} \in {\Bbb M}({\cal N},


E)$ и $A \in {{\cal P}(E)}$,


то $(A_i)_{i \in {\cal N}} \downarrow A$ означает:


1)~$A$ есть пересечение всех множеств ${A_i}$, $i \in {\cal N}$;


2)~${A_{k+1}} \subset {A_k}$


$\fo{k} \in {\cal N}$.


Если же $A$ и $B$ --- множества, $({\alpha}_i)_{i \in {\cal N}} \in


{{\Bbb M}(A,B)}^{\cal N}$ и $\alpha \in {\Bbb M}(A,B)$,


то $\Def$


$$


({({\alpha}_i)_{i \in {\cal N}}} \Downarrow \alpha) \Longleftrightarrow


(\fo{a} \in A:{({{\alpha}_i}(a))_{i \in {\cal N}}} \downarrow


{{\alpha}(a)}).


$$


Итак, введены два типа сходимости, один из которых


реализуется в пространстве п/м того или иного множества,


а другой --- в пространстве МО. В пространствах первого типа


порядок определяется вложениями, а в пространствах второго типа


используется соглашение: если


$A$, $B$ --- множества, $p \in {\Bbb M}(A,B)$ и $q \in


{\Bbb M}(A,B)$, то $\Def:


(p \sqsubseteq q) \Longleftrightarrow (\fo{a} \in A:


{p(a)} \subset {q(a)})$;


имеем порядок в ${\Bbb M}(A,B)$. Если ${\cal S}$ ---


непустое семейство, то


${{\cal S}^{\cal S}_{(0)}} \stac \{\alpha \in {\cal S}^{\cal S} \mid


\fo{E} \in {\cal S}:{{\alpha}(E)} \subset E\}$.


Если $H$ --- множество, то ([19], c.\,321) $(\sigma -{\downarrow})[H]$ есть


$\Def$ множество всех ${\cal H} \in {{\cal P}'({\cal P}(H))}$


таких, что


$\fo{(L_i)_{i \in {\cal N}}} \in {{\cal H}^{\cal N}}$ $\fo{L} \in


{{\cal P}(H)}:


({(L_i)_{i \in {\cal N}}} \downarrow L) \Lo (L \in {\cal H})$;


элементы $(\sigma - {\downarrow})[H]$ суть непустые секвенциально


$\downarrow$-замкнутые семейства п/м $H$ и только они. Для всяких


множества $H$ и семейства ${\cal H} \in (\sigma - {\downarrow})[H]$


введем множество ([19], c.\,321)


${{\frak A}_H}({\cal H})$,


${{\frak A}_H}({\cal H}) \subset {{\cal P}(H)}^{{\cal P}(H)}_{(0)}$;


в терминах ${{\frak A}_H}({\cal H})$


реализуется ([19], c.\,322) свойство экстремальной по порядку


неподвижной точки, восходящее к [8], [9], [15].


Если $A$ и $B$ --- множества, то


${\cal M}[A;B] \stac \{\alpha \in {\Bbb M}(A,B)^{{\Bbb M}(A,B)} \mid


\fo{\cal H} \in {\Bbb M}(A,B):{{\alpha}({\cal H})} \sqsubseteq


{\cal H}\}$;


при $\beta \in {\cal M}[A;B]$ и ${\cal C} \in {\Bbb M}(A,B)$


имеем сходимость ${({{\beta}^k}({\cal C}))_{k \in {\cal N}}}


\Downarrow {{\beta}^\infty}({\cal C})$. Аналогично, если


${\cal S}$ есть семейство всех п/м некоторого множества,


$\alpha \in {\cal S}^{\cal S}_{(0)}$ и $E \in {\cal S}$,


то ${({{\alpha}^k}(E))_{k \in {\cal N}}}


\downarrow \up{\infty}{\alpha}(E)$.


Если $({\Bbb T},\tau)$ --- ТП и


$x \in {\Bbb T}$, то ${N_\tau}(x)$ есть $\Def$ фильтр


всех окрестностей ([25], c.\,62)


$x$ в $({\Bbb T},\tau)$ ([26], c.\,100). Используем сходимость по


Мору--Смиту ([25], гл.\,2). Для обозначения направленности в


множестве $H$ используем триплет вида $(D,\preceq,f)$:


$(D,\preceq)$, $D \ne \emp$, --- НМ и


$f \in {H^D}$;


$(H - {\ass})[D;\preceq;f] \stac \{S \in {{\cal P}(H)} \mid


\exi{d} \in D$ $\fo{\delta} \in D:(d \preceq \delta) \Lo


(f(\delta) \in S)\}$


есть фильтр, ассоциированный с $(D,\preceq,f)$. Если $({\Bbb T},\tau)$


есть ТП, $(D,\preceq,f)$ --- направленность в ${\Bbb T}$ и $x \in


{\Bbb T}$, то, как обычно, полагаем $\Def$


$$


((D,\preceq,f) @>{\tau\,}>> x) \Longleftrightarrow ({N_\tau}(x)


\subset ({\Bbb T} - {\ass})[D;\preceq;f]).


$$


Если $(D,\preceq)$ и $(\Delta,\ll)$ ---


непустые НМ, то $(\Isot)[D;\preceq;\Delta;\ll]$ соответствует


([23], c.\,68) (множество всех $(\preceq,\ll)$-изотонных отображений


из множества ${{\Delta}^D}$ с конфинальным образом $D$).


Если $(D,\preceq)$ --- непустое НМ,


то полагаем $(\isot)[D;\preceq] \stac (\Isot)[D;\preceq;{\cal N};\le]$,


где $\le$ --- естественный порядок ${\cal N}$.


Используем характеризацию компактности


ТП в терминах изотонного прореживания ([23], c.\,69) направленностей


до сходящихся поднаправленностей ([25], гл.\,2).





\section{Обобщенная динамика}





Пусть ${t_0} \in {\Bbb R}$, ${{\vt}_0} \in\, ]{t_0},\infty[$ и ${I_0}


\stac [{t_0},{{\vt}_0}]$. Фиксируем непустое множество $X$, а также


множество ${\Bbb C} \in {{\cal P}'}(X^{I_0})$; $X$~играет роль фазового


пространства системы, ${\Bbb C}$ определяет запас траекторий в


$X$. Пусть ${\cal T}$ --- топология множества ${\Bbb C}$,


фиксированная в дальнейшем, а ${\Bbb F}$ и ${\frak K}$ --- соответственно


семейства всех замкнутых и всех компактных ([27], c.\,196) в ТП


$({\Bbb C},{\cal T})$ п/м ${\Bbb C}$. Используем ${\bold D} \stac


{I_0} \times {\Bbb C}$ в качестве множества позиций;


для $z \in {\bold D}$ элементы ${{\pr_1}(z)} \in {I_0}$ и ${{\pr_2}(z)} \in


{\Bbb C}$ (компоненты $z$)


$\Def$ таковы, что $z = ({{\pr}_1}(z),{{\pr_2}(z)})$.


Пусть $\Upsilon$ --- непустое множество и $\Om \in


{{\cal P}'}({\Upsilon}^{I_0})$; рассмотрим систему


\begin{equation}


S:{\bold D} \times \Om \lo {{\cal P}'}({\Bbb C});


\end{equation}


$S(z,\omega)$ имеет смысл


пучка траекторий ($\omega \in \Om$ играет


роль реализации неопределенных факторов, а $z \in {\bold D}$ --- роль


начальной позиции).





\begin{notes}


В (3.1) при формировании ``позиции'' $z = (t,g)$ используется функция


$g \in {\Bbb C}$, определенная на ${I_0}$, а не на $[{t_0},t]$ (если только


$t \ne {{\vt}_0}$). В то же время в конкретных задачах управления


используется, как правило, лишь $(g \mid [{t_0},t])$, т.е. имеет


место наследственность реакций в (3.1).


\end{notes}





В частном случае значение (3.1) определяется при


${t_{00}} \in {I_0}$ в виде пучка, получаемого склеиванием решений


векторного дифференциального уравнения


\begin{equation}


\Dot x = f(t,x,v) + {B(t,x)}u,\quad u \in P,\ \;v \in Q,


\end{equation}


где $t \in [{t_{00}},{{\vt}_0}]$, с начальной историей, определенной на


$[{t_0},{t_{00}}]$.


Постулируем в


отношении (3.2) традиционные условия [2], [3]: непрерывность правой части


(3.2) по совокупности переменных, локальная липшицевость и условие


подлинейного роста по $x$.


Полагаем, что $P$ и $Q$ ---


непустые компакты в конечномерных арифметических пространствах,


$P$ --- выпуклое множество.


В качестве фазового


пространства системы (3.2) используем ${{\Bbb R}^n} = X$;


через ${C_n}(I_0)$


обозначаем множество всех непрерывных $n$-вектор-функций на ${I_0}$.


Решения (3.2), понимаемые в смысле Каратеодори [28], рассматриваем на


$I_0$ или на $[t,{{\vt}_0}]$, где $t \in {I_0}$.


Полагаем, что ${\cal U}$ и ${\cal V} = \Om$ суть множества всевозможных


борелевских функций из $I_0$ в $P$ и в $Q$ соответственно.


Если ${t_*} \in {I_0}$, ${x_*} \in


{{\Bbb R}^n}$, $U \in {\cal U}$ и $V \in {\cal V}$, то


через ${\vp}(\cdot,{t_*},{x_*},U,V) = ({\vp}(t,{t_*},{x_*},U,V) \in


{{\Bbb R}^n}$, ${t_*} \le t \le {{\vt}_0})$ обозначаем решение


(3.2) в смысле Каратеодори при подстановке $u = {U(t)}$, $v = {V(t)}$ и при


начальных условиях ${\vp}({t_*},{t_*},{x_*},U,V) = {x_*}$. Определяем


${\Bbb C}$ в виде ${C_n}(I_0)$. Если ${t_*} \in {I_0}$,


$g \in {C_n}(I_0)$,


$U \in {\cal U}$ и $V \in {\cal V}$, то ${\wt \vp}(\cdot,{t_*},g,U,V)$


определяем склеиванием $g$ и ${\vp}(\cdot,{t_*},{g(t_*)},U,V)$ в момент


${t_*}$.


При ${t_*} \in {I_0}$, $g \in {\Bbb C}$ и


$\omega = V \in {\cal V}$ пучок всех траекторий ${\wt \vp}(\cdot,


{t_*},g,U,\omega)$, $U \in {\cal U}$, рассматриваем как


значение $S$ (3.1), т.\,е. как множество $S(({t_*},g),\omega)$.


Отображение (3.1) в нашем случае наследственно и компактнозначно в смысле


топологии равномерной сходимости пространства ${\Bbb C}$;


о других известных свойствах см. [28], [29]. Обсудим совсем кратко


другой вариант (3.1), обращаясь к ([20]; [30], c.\,136). Рассмотрим сначала


систему


\begin{equation}


\Dot x (t) = {A(t)}{x(t)} + {f(t)}{b(t)} +


{C(t)}{{\omega}(t)}


\end{equation}


в $n$-мерном фазовом пространстве; здесь $t


\in [a,{{\vt}_0}]$, где $a \in {I_0}$. Ограничения на выбор


$\omega$ полагаем прежними (${{\omega}(t)} \in Q$, где


$Q$ --- конечномерный компакт); сохраняем прежнее предположение


относительно $\Om$. Полагаем, что $A(\cdot)$ и $C(\cdot)$ суть


непрерывные на $I_0$ матрицанты, а вектор-функция $b(\cdot)$ определена


на $I \stac [{t_0},{{\vt}_0}[$ и является равномерным пределом


последовательности кусочно-постоянных (к.-п.) и непрерывных


справа (н.\,спр.) $n$-вектор-функций на $I$. В качестве $f$ допускаем (в


(3.3)) использование произвольной к.-п. и н.\,спр. вещественнозначной (в/з)


функции на $I$, для которой


${\int\limits_{t_0}^{{\vt}_0}}|{f(t)}|{dt}


\le c$, где $c \in [0,\infty[$. Разумеется, в (3.3) используется лишь


$(f \mid [a,{{\vt}_0}[)$ и допускается , что часть энергоресурса


истрачена до момента $a$, соответствующего текущей


задаче управления системой (3.3). В этой связи полагаем, что


фазовый вектор системы дополняется особой ресурсной компонентой, получаемой


в момент $\vt \in {I_0}$


интегрированием функции $t \longmapsto |f(t)|$ на


$[{t_0},\vt[$, а ресурс управления системой (3.3) из начального состояния,


сложившегося к моменту $a$, определяется как разность $c$ и


интеграла начального отрезка функции $|f|$.


При этом в качестве основной используется


задача управления на отрезке $I_0$ из начальной позиции


$({t_0},{x_0})$, где $x_0$ --- заданный $n$-мерный вектор.


Ограничимся данным


обсуждением этой известной (см., напр., [2])


конструкции и отметим, что по ряду причин пространство


упомянутых простейших управлений требует расширения, которое в ([20];


[30], c.\,136) реализовано в классе конечно-аддитивных (к.-а.) мер


ограниченной вариации


при использовании


обобщенного аналога формулы Коши. В итоге возникают разрывные, вообще говоря,


траектории. Разумеется, обобщенная формула Коши ([30], c.\,136) должна быть


дополнена соотношением, определяющим изменение вариации к.-а. (обобщенного)


управления-меры с течением времени; получаем при этом расширение


исходной ресурсной компоненты пополненной версии системы (3.3).


Множество обычных (к.-п. и н.\,спр.) управлений $f$ вкладывается при этом


в виде всюду плотного множества в компакт, определяемый в виде шара


радиуса $c$ в пространстве в/з к.-а. мер ограниченной вариации с областью


определения в виде полуалгебры ([30], cc.\,57,~89) всех промежутков


$[{\alpha},{\beta}[$, ${t_0} \le \alpha \le \beta \le {{\vt}_0}$.


По аналогии со


случаем системы (3.2) введем склеивание с


начальными историями, если в расширенной версии (3.3)


$a \ne {t_0}$. В качестве ${\Bbb C}$ здесь


можно


использовать множество всех функций, действующих


из ${I_0}$ в арифметическое


пространство размерности $n + 1$, у каждой из которых


последняя (для определенности)


компонента принимает значения из $[0,c]$, не возрастает и в момент


$t_0$ принимает значение $c$;


соответствующая топология


поточечной сходимости может использоваться в качестве ${\cal T}$.





Возвращаясь к общей постановке, рассмотрим серию условий.





\begin{conds}


$S(z,\omega) \in {\Bbb F}$


$\fo{z} \in {\bold D}$, $\fo{\omega} \in \Om$.


\end{conds}





\begin{conds}


Для всяких $t \in {I_0}$ и $\omega \in \Om$ множество $\{(y,h) \in


{\Bbb C} \times {\Bbb C} \mid h \in S((t,y),\omega)\}$ замкнуто в ${\Bbb


C} \times {\Bbb C}$ с топологией произведения ([25], c.\,125)


двух экземпляров ТП $({\Bbb C},{\cal T})$.


\end{conds}





\begin{conds}


$\fo{t_*} \in {I_0}$, $\fo{u_*} \in {\Bbb C}$,


$\fo{\omega} \in \Om$ $\exi{H} \in


{N_{\cal T}}(u_*)$,


$\exi{K} \in {\frak K}$ $\fo{u} \in H:S(({t_*},u),\omega)


\subset K$.


\end{conds}





\begin{conds}


$h \in S((t,h),\omega)$


$\fo{z} \in {\bold D}$, $\fo{\omega} \in \Om$, $\fo{h} \in S(z,\omega)$,


$\fo{t} \in[{\pr_1}(z),{{\vt}_0}[$.


\end{conds}





\begin{conds}


$((u \mid [{t_0},t]) = (v \mid [{t_0},t])) \Lo (\fo{\omega} \in \Om:


S((t,u),\omega) = S((t,v),\omega))$


$\fo{t} \in {I_0}$, $\fo{u} \in {\Bbb C}$, $\fo{v} \in {\Bbb C}$.


\end{conds}





Условие 3.1 следует из условия 3.2. При условиях 3.1,


3.3 $\fo{z} \in {\bold D}$, $\fo{\omega} \in \Om:S(z,\omega) \in


{\frak K}$. Условие 3.5 означает наследственность $S$ (3.1); если к тому


же верно условие 3.4, то справедливо полугрупповое свойство.





\section{Оператор программного поглощения\protect\\


и непрямая итерационная процедура}





Рассмотрим процедуру, подобную [8], [9].


В ее основе --- ОПП, определяемый аналогично [9].


Если $H \in {{\cal P}({\bold D})}$ и $t \in {I_0}$, то ${H\langle{t}\rangle}


\stac \{g \in {\Bbb C} \mid (t,g) \in H\}$. Полагаем


\begin{equation}


{\bold F} \stac \{F \in {{\cal P}({\bold D})} \mid \fo{t} \in {I_0}:


F\langle{t}\rangle \in {\Bbb F}\}.


\end{equation}


Получили семейство всех п/м ${\bold D}$, замкнутых в


топологии произведения ([25], c.\,125) ТП $({I_0},


{{\cal P}(I_0)})$ и $({\Bbb C},{\cal T})$; $({I_0},{{\cal P}(I_0)})$


соответствует оснащению $I_0$ дискретной топологией ([25], c.\,61);


${\bold F} \in (\sigma - {\downarrow})[{\bold D}]$.


Фиксируем множество $M \in {\Bbb F}$. Осуществление траекторий


в пределах $M$ рассматриваем как целевую установку при выборе


(конкретной траектории) из


пучка, реализующегося в силу (3.1).


На промежуточных


этапах решения используем программные (по смыслу) задачи с функциональными


ограничениями. В этой связи полагаем $\fo{H} \in {{\cal P}({\bold D})}$,


$\fo{z} \in {\bold D}$, $\fo{\omega} \in \Om$


\begin{equation}


{\Pi}(\omega \mid z,H) \stac \{h \in S(z,\omega) \cap M \mid \fo{t} \in


[{\pr_1}(z),{{\vt}_0}[\,:(t,h) \in H\}.


\end{equation}


При условии 3.1 из (4.2) имеем


$\fo{F} \in {\bold F}$, $\fo{z} \in {\bold D}$,


$\fo{\omega} \in \Om:{\Pi}(\omega \mid z,F) \in {\Bbb F}$.





Введем ОПП ${\Bbb A}:{{\cal P}({\bold D})} \lo {{\cal P}({\bold D})}$,


полагая $\fo{H} \in {{\cal P}({\bold D})}$


\begin{equation}


{{\Bbb A}(H)} \stac \{z \in H \mid \fo{\omega} \in \Om:


{\Pi}(\omega \mid z,H) \ne \emp\};


\end{equation}


из (4.3) видно, что ${\Bbb A} \in {{\cal P}({\bold D})}^{{\cal P}({\bold


D})}_{(0)}$; ОПП (4.3) является развитием аналогичной конструкции


[8], [9]. Естественно привлечь для описания ОПП неигровые операторы


подобно [12]: если $\omega \in \Om$, то


${{\bold A}_\omega} \in {{\cal P}({\bold D})}^{{\cal P}({\bold D})}_{(0)}$


есть $\Def$ оператор, для которого $\fo{H} \in {{\cal P}({\bold D})}:


{{\bold A}_\omega}(H)\stac\{z\in H \mid {\Pi}(\omega \mid z,H) \ne \emp\}$.


В терминах $({\bold A}_\omega)_{\omega \in \Om}$ получаем очевидное


представление ОПП ${\Bbb A}$. Именно, имеем свойство (см.(4.3)): если


$U \in {{\cal P}({\bold D})}$, то ${\Bbb A}(U)$ есть пересечение всех


множеств ${{\bold A}_\omega}(U)$, $\omega \in \Om$.





Очевидным следствием определений является


при условии 3.4 равенство


${{\bold A}_\omega}


= {{\bold A}_\omega} \circ {{\bold A}_\omega}$


$\fo{\omega} \in \Om$. В итоге справедливо





\begin{props}


При условии $3.4$ 


$$


\{H \in {{\cal P}({\bold D})} \mid


H = {{\bold A}_\omega}(H)\} = \{{{\bold A}_\omega}(H):H \in


{{\cal P}({\bold D})}\}\ \; \fo{\omega} \in \Om.


$$


\end{props}





В общем случае имеет место свойство


$$


(H = {\Bbb A}(H)) \Longleftrightarrow


(\fo{\omega} \in \Om:H = {{\bold A}_\omega}(H)) 


\ \, \fo{H} \in {{\cal P}({\bold D})}.


$$





\begin{props}


При условиях $3.2$, $3.3$


${{\bold A}_\omega}(F) \in {\bold F}$


$\fo{\omega} \in \Om$, $\fo{F} \in {\bold F}$.


\end{props}





\begin{pf}


Пусть выполнены условия 3.2, 3.3, $\omega \in \Om$,


$F \in {\bold F}$ и ${t_*} \in {I_0}$. Рассмотрим


${{\bold A}_\omega}(F)\langle{t_*}\rangle \in {{\cal P}({\Bbb C})}$. Пусть


\begin{equation}


g \in \mathrm{cl}({{\bold A}_\omega}(F)\langle{t_*}\rangle,{\cal T}).


\end{equation}


В силу теоремы Биркгофа [25]--[27] для $g \in {\Bbb C}$ можно указать


направленность $({\cal D},\preceq,\vp)$ в


${{\bold A}_\omega}(F)\langle{t_*}\rangle$ со свойством


\begin{equation}


({\cal D},\preceq,\vp) @>{\cal T\,}>> g.


\end{equation}


При $d \in {\cal D}$ полагаем ${z_d^*} \stac ({t_*},{{\vp}(d)})$;


${z_d^*} \in {{\bold A}_\omega}(F)$. Из (4.5) в силу


${F\langle{t_*}\rangle} \in {\Bbb F}$ имеем $g \in {F\langle{t_*}\rangle}$,


т.\,е. ${z_*} \stac ({t_*},g) \in F$. Рассмотрим отображение


$d \longmapsto {\Pi}(\omega \mid {z_d^*},F):{\cal D} \lo


{{\cal P}'}(M)$.


Пусть $\rho \in {\prod\limits_{d \in {\cal D}}}{\Pi}(\omega \mid


{z_d^*},F)$. С учетом условия 3.3 подберем такие


$H \in {N_{\cal T}}(g)$ и $K \in {\frak K}$, что $\fo{u} \in H:


S(({t_*},u),\omega) \subset K$. Тогда $H \in ({\Bbb C} - {\ass})[{\cal


D};\preceq;\vp]$. Подберем ${\bold d} \in {\cal D}$ такое, что


$\fo{d} \in {\cal D}:({\bold d} \preceq d) \Lo ({\vp}(d) \in H)$.


Введем ${{\cal D}_0} \stac \{d \in {\cal D} \mid {\bold d} \preceq d\}$,


получая $\preceq$-конфинальное ([5], гл.\,2) п/м ${\cal D}$.


Оснащаем ${{\cal D}_0} \in {{\cal P}'}({\cal D})$


направлением $\sqsubseteq\, \stac\, \preceq \cap ({{\cal D}_0} \times


{{\cal D}_0})$, для которого при ${d_1} \in {{\cal D}_0}$ и


${d_2} \in {{\cal D}_0}$ эквивалентны свойства


1)~${d_1} \sqsubseteq {d_2}$; 2)~${d_1} \preceq {d_2}$. Пусть


$\overline{\vp} \stac (\vp \mid {{\cal D}_0})$; тогда


$({{\cal D}_0},\sqsubseteq,\overline{\vp})$ --- направленность в


$H$. Для $\overline{\rho} \stac (\rho \mid {{\cal D}_0})$


имеем: $({{\cal D}_0},\sqsubseteq,\overline{\rho})$ есть направленность в


$K$. С учетом компактности $K$ подберем $h \in K$, НМ


$(\Delta,\ll)$, $\Delta \ne \emp$,


и $l \in (\Isot)[\Delta;\ll;{{\cal D}_0};\sqsubseteq]$,


для которых


\begin{equation}


(\Delta,\ll,{\overline{\rho}} \circ l) @>{\cal T\,}>> h.


\end{equation}


Разумеется, $(\ov{\rho} \circ l)(\delta) \in {\Pi}(\omega \mid


{z_{l(\delta)}^*},F)$ при $\delta \in \Delta$. В силу $F \in {\bold F}$


имеем из (4.6) свойство $(t,h) \in F$ $\fo{t} \in [{t_*},{{\vt}_0}[$.


Кроме того, из (4.6) следует, что $h \in M$.


Заметим, что $(\ov{\rho} \circ l)(\delta) \in S({z_{l(\delta)}^*},\omega)$


при $\delta \in \Delta$. С другой стороны, в силу условия 3.2


множество


$\Xi \stac \{(y,s) \in {\Bbb C} \times {\Bbb C} \mid s \in S(({t_*},y),


\omega)\} \in {\cal P}({\Bbb C} \times {\Bbb C})$


замкнуто в ${\Bbb C} \times {\Bbb C}$ с топологией произведения двух


экземпляров ТП $({\Bbb C},{\cal T})$. Из (4.5) следует сходимость


$(\Delta,\ll,\ov{\vp} \circ l)$ к $g$ в ТП $({\Bbb C},{\cal T})$.


С учетом (4.6) и определения ${z_d^*}$, $d \in {\cal D}$, имеем свойства


$$


((\Delta,\ll,\ov{\vp} \circ l) @>{\cal T\,}>> g) \&


((\Delta,\ll,\ov{\rho} \circ l) @>{\cal T\,}>> h) \&


(\fo{\delta} \in \Delta:((\ov{\vp} \circ l)(\delta),(\ov{\rho} \circ


l)(\delta)) \in \Xi).


$$


Как следствие, $(g,h) \in \Xi$, т.\,е. $h \in S(({t_*},g),\omega)$.


С учетом (4.2) имеем $h \in {\Pi}(\omega \mid {z_*},F)$. Итак,


${z_*} \in F$ обладает свойством ${\Pi}(\omega \mid {z_*},F) \ne


\emp$, т.\,е. ${z_*} \in {{\bold A}_\omega}(F)$. Как следствие,


$g \in {{\bold A}_\omega}(F)\langle{t_*}\rangle$. Поскольку выбор


(4.4) был произвольным, свойство


${{\bold A}_\omega}(F)\langle{t_*}\rangle \in


{\Bbb F}$ установлено. Но и выбор $t_*$ также был произвольным, что


означает ${{\bold A}_\omega}(F) \in {\bold F}$.


\end{pf}





Очевидным является





\begin{props}


Если выполнено условие $3.4$, то имеет место равенство


$\{H \in {{\cal P}({\bold D})} \mid H = {\Bbb A}(H)\} =


{\bigcap\limits_{\omega \in \Om}}\{{{\bold A}_\omega}(L):L \in


{{\cal P}({\bold D})}\}$.


\end{props}





Из предложений 4.2, 4.3 следует





\begin{props}


Если выполнены условия $3.2$--$3.4$, то справедливо равенство


$\{F \in {\bold F} \mid F = {\Bbb A}(F)\} = {\bigcap\limits_{\omega \in


\Om}}\{{{\bold A}_\omega}(F):F \in {\bold F}\}$.


\end{props}





\begin{props}


При условиях $3.2$, $3.3$


${\Bbb A}(F) \in {\bold F}$\;


$\fo{F} \in {\bold F}$.


\end{props}





Доказательство непосредственно следует из (4.1) и предложения 4.2.





\begin{props}


При условиях $3.1$, $3.3$ ОПП секвенциально непрерывен на ${\bold F}{:}$


$\fo{(F_i)_{i \in {\cal N}}} \in {{\bold F}^{\cal N}}$, $\fo{F} \in


{{\cal P}({\bold D})}$


\begin{equation}


({(F_i)_{i \in {\cal N}}} \downarrow F) \Lo


({({\Bbb A}(F_i))_{i \in {\cal N}}} \downarrow {{\Bbb A}(F)}).


\end{equation}


\end{props}





\begin{pf}


Пусть $(F_i)_{i \in {\cal N}} \in {{\bold F}^{\cal


N}}$ и $F \in {{\cal P}({\bold D})}$ таковы, что истинна посылка в (4.7).


Как следствие, ${\Bbb A}(F)$ есть п/м пересечения всех множеств


${{\Bbb A}(F_k)}$, $k \in {\cal N}$; кроме того,


${{\Bbb A}(F_{i+1})} \subset {{\Bbb A}(F_i)}$ $\fo{i} \in {\cal N}$.


Выберем произвольно ${z_*} \in


{\bigcap\limits_{k \in {\cal N}}}{{\Bbb A}(F_k)}$.


Разумеется, ${z_*} \in F$. Пусть $\omega \in \Om$; тогда


$({\Pi}(\omega \mid {z_*},{F_k}))_{k \in {\cal N}}$ ---


последовательность в ${{\cal P}'}({\Bbb C})$. Пусть


$(h_k)_{k \in {\cal N}}$ --- элемент произведения всех множеств


${\Pi}(\omega \mid {z_*},{F_k})$, $k \in {\cal N}$;


$(h_k)_{k \in {\cal N}}$ --- последовательность в


$S({z_*},\omega) \cap M$, где


$S({z_*},\omega) \in {\frak K}$. В частности, множество


$S({z_*},\omega)$ счетно-компактно в $({\Bbb C},{\cal T})$.


Поэтому ([23], c.\,69) можно указать $h \in S({z_*},\omega)$, непустое НМ


$({\cal D},\preceq)$, а также $l \in (\isot)[{\cal D};\preceq]$,


для которых


\begin{equation}


({\cal D},\preceq,(h_{l(d)})_{d \in {\cal D}}) @>{\cal T\,}>> h.


\end{equation}


Тогда $h \in S({z_*},\omega) \cap M$. При ${t^*} \in


[{\pr_1}({z_*}),{{\vt}_0}[$


имеем


\begin{equation}


{(h_k)_{k \in {\cal N}}} \in {\prod\limits_{k \in


{\cal N}}}{F_k}\langle{t^*}\rangle.


\end{equation}


Фиксируя $n \in {\cal N}$, получаем ${F_n}\langle{t^*}\rangle \in


{\Bbb F}$ и в силу (4.8), (4.9) $h \in {F_n}\langle{t^*}\rangle$


(см. рассуждение [23], c.\,71). Поскольку выбор $n$ был


произвольным, то $({t^*},h) \in F$. Однако и $t^*$ выбиралось


произвольно. Поэтому $h \in {\Pi}(\omega \mid {z_*},F)$.


Установлено $\fo{\omega} \in \Om:{\Pi}(\omega \mid {z_*},F) \ne


\emp$; итак, ${z_*} \in {{\Bbb A}(F)}$. Этим завершается


обоснование (4.7).


\end{pf}





\begin{props}


При условиях $3.2$, $3.3$ имеет место


${\Bbb A} \in {{\frak A}_{\bold D}}({\bold F})$.


\end{props}





Доказательство получается комбинацией (4.3), предложений


4.5 и 4.6.





\begin{thms}


При условиях $3.2$, $3.3$ множество $\up{\infty}{\Bbb A}({\bold D})$


есть экстремальная по порядку неподвижная точка ОПП ${\Bbb A}:$


$$


({\Bbb A}(\up{\infty}{\Bbb A}({\bold D})) =


\up{\infty}{\Bbb A}({\bold D})) \& (\fo{H} \in


{\cal P}({\bold D}):(H = {{\Bbb A}(H)}) \Lo (H \subset


\up{\infty}{\Bbb A}({\bold D}))) \&


({({{\Bbb A}^k}({\bold D}))_{k \in {\cal N}}} \downarrow


\up{\infty}{\Bbb A}({\bold D})).


$$


\end{thms}





Для доказательства достаточно (наряду с предложением 4.7) учесть


предложение 2.1 работы [19] (см. подобные конструкции в [8], [12], [15]).


Разумеется, теорема 4.1 означает


сходимость последовательности итераций на основе ОПП ${\Bbb A}$ (с начальным


элементом ${\bold D}$) к наибольшей неподвижной точке ОПП. Заметим, что


$\up{\infty}{\Bbb A}({\bold D}) \in {\bold F}$. Отметим


свойства, связанные с динамикой. Пусть ${{\cal P}_{\bold h}}({\bold D})$


есть $\Def$ семейство всех множеств $H \in {{\cal P}({\bold D})}$ таких, что


$\fo{t} \in {I_0}$, $\fo{u} \in {\Bbb C}$, $\fo{v} \in {\Bbb C}$


\begin{equation}


(((t,u) \in H) \& ((u \mid [{t_0},t]) = (v \mid [{t_0},t]))) \Lo


((t,v) \in H).


\end{equation}


Свойство (4.10) имеет смысл наследственности и естественно


в условиях, когда (3.1) определяет динамическую систему.


Если ${\cal H} \in {{\cal P}'}({{\cal P}_{\bold h}}({\bold D}))$,


то пересечение всех множеств из ${\cal H}$ само является элементом


${{\cal P}_{\bold h}}({\bold D})$,


т.\,е. наследственным множеством. Очевидным является





\begin{props}


При условии $3.5$\;


$\fo{t} \in {I_0}$, $\fo{u} \in {\Bbb C}$,


$\fo{v} \in {\Bbb C}:


((u \mid [{t_0},t]) = (v \mid [{t_0},t])) \Lo (\fo{H} \in {{\cal


P}({\bold D})}$, $\fo{\omega} \in \Om:{\Pi}(\omega \mid (t,u),H) =


{\Pi}(\omega \mid (t,v),H))$.


\end{props}





Отсюда и из определений вытекает





\begin{cors}


При условии 3.5


${{\bold A}_\omega}(H) \in {{\cal P}_{\bold h}}({\bold


D})$ $\fo{H} \in {{\cal P}_{\bold h}}({\bold D})$,


$\fo{\omega} \in \Om$.


\end{cors}





Из следствия 4.1


и соотношений для ${\Bbb A}$ и ${{\bold A}_\omega},\;\omega \in \Om$,


получаем





\begin{props}


При условии $3.5$


${\Bbb A}(H) \in {{\cal P}_{\bold h}}({\bold D})$\;


$\fo{H} \in {{\cal P}_{\bold h}}({\bold D})$.


\end{props}





Поскольку ${\bold D} \in {{\cal P}_{\bold h}}({\bold D})$, то


из предложения 4.9


имеем при условии 3.5 свойство: $({{\Bbb A}^k}({\bold D}))_{k \in


{{\cal N}_0}}$ есть последовательность в ${{\cal P}_{\bold h}}({\bold D})$.


Из ранее упомянутого свойства полной мультипликативности


${{\cal P}_{\bold h}}({\bold D})$ получаем при условии 3.5 свойство


$\up{\infty}{\Bbb A}({\bold D}) \in {{\cal P}_{\bold h}}({\bold D})$. В


дальнейшем рассмотрим один ослабленный вариант условия наследственности.





\section{Итерации в пространстве многозначных отображений}





Рассмотрим вариант конструкции [20]--[23], [31], которая в разделе 1


была названа прямой, поскольку при ее реализации


в задачах теории ДИ осуществляется непосредственный перевод


многозначного аналога псевдостратегии [32] в многозначную


квазистратегию ([7]--[10], [12]--[14], [33], [34]; упомянутый тип управляющих


процедур можно рассматривать как многозначный аналог квазистратегий


[35]). Отметим, что основное отличие квазистратегий от псевдостратегий


связано с постулируемым в первом случае условием неупреждаемости


(физической осуществимости) отклика на реализации неопределенных


факторов; псевдостратегии таким свойством могут не обладать.





Фиксируем ${{\tau}_0} \in [{t_0},{{\vt}_0}[$ и ${g_0} \in {\Bbb C}$;


пару ${z_0} \stac ({{\tau}_0},{g_0}) \in {\bold D}$ рассматриваем далее


в качестве фиксированной начальной позиции, а


\begin{equation}


{{\cal C}_0} \stac {\Pi}(\cdot \mid {z_0},{\bold D}) =


({\Pi}(\omega \mid {z_0},{\bold D}))_{\omega \in \Om}


\end{equation}


используем в качестве целевого МО ([20]--[23], [31]);


в задачах теории ДИ можно интерпретировать (5.1) как


многозначную псевдостратегию. В силу (4.2) имеем для


${{\cal C}_0} \in {\Bbb M}(\Om,{\Bbb C})$ свойство:


${{\cal C}_0}(\omega) = S({z_0},\omega) \cap M$ при $\omega \in \Om$.


Пусть $\fo{t} \in [{{\tau}_0},{{\vt}_0}]$


\begin{multline}


(\fo{\omega} \in \Om:{{\wt \Om}_t}(\omega) \stac \{{\wt \omega} \in


\Om \mid (\omega \mid [{t_0},t]) = ({\wt \omega} \mid [{t_0},t])\})\&\\


%


\& (\fo{h} \in {\Bbb C}:{{\wt Z}_t}(h)


\stac \{{\wt h} \in {\Bbb C}


\mid (h \mid [{t_0},t]) = ({\wt h} \mid [{t_0},t])\}).


\end{multline}


В терминах (5.2) вводим оператор $\Gm \in {\cal M}[\Om;{\Bbb C}]$:


$\fo{\cal C} \in {\Bbb M}(\Om,{\Bbb C})$, $\fo{\omega} \in \Om$


\begin{equation}


{\Gm}({\cal C})(\omega) \stac \{h \in {{\cal C}(\omega)} \mid \fo{t} \in


[{{\tau}_0},{{\vt}_0}]\ \fo{\wt \omega} \in {{\wt \Om}_t}(\omega):


{{\wt Z}_t}(h) \cap {{\cal C}({\wt \omega})} \ne \emp\}.


\end{equation}


Оператор $\Gm$ (5.3) можно рассматривать в виде конкретизации


аналогичного оператора [20]--[23], [31]. Пусть


\begin{equation}


{\bold N} \stac \{{\cal C} \in {\Bbb M}(\Om,{\Bbb C}) \mid ({\cal C} =


{\Gm}({\cal C})) \& ({\cal C} \sqsubseteq {{\cal C}_0})\}


\end{equation}


(множество всех неупреждающих м/с ${\cal C}_0$); ${\bold N} \ne \emp$.


По аналогии с [22], [31] определяется


$\sqsubseteq$-наибольший элемент ${\bold N}$ (5.4) в виде


МО


\begin{equation}


(\na)[{\cal C}_0] \stac \Big({\bigcup\limits_{{\cal C} \in


{\bold N}}}{{\cal C}(\omega)}\Big)_{\omega \in \Om} \in {\bold N};


\end{equation}


ясно, что ${\cal L} \sqsubseteq (\na)[{\cal C}_0]$ при ${\cal L} \in


{\bold N}$. С $\Gm$ связываем последовательность


${({\Gm}^k)_{k \in {{\cal N}_0}}}:{{\cal N}_0} \lo


{{\Bbb M}(\Om,{\Bbb C})}^{{\Bbb M}(\Om,{\Bbb C})}$


и предельный оператор ${{\Gm}^\infty}$, действующий в


${\Bbb M}(\Om,{\Bbb C})$.


Для нас существенно построение последовательности


$({{\Gm}^k}({\cal C}_0))_{k \in {{\cal N}_0}}$ в ${\Bbb M}(\Om,{\Bbb C})$


и МО ${{\Gm}^\infty}({\cal C}_0)$;


\begin{equation}


{({{\Gm}^k}({\cal C}_0))_{k \in {\cal N}}} \Downarrow


{{\Gm}^\infty}({\cal C}_0).


\end{equation}


Разумеется, (5.6) реализует предел


последовательности итераций


\begin{equation}


({{\Gm}^0}({\cal C}_0) = {{\cal C}_0}) \& (\fo{k} \in {\cal N}:


{{\Gm}^k}({\cal C}_0) = {\Gm}({{\Gm}^{k-1}}({\cal C}_0))).


\end{equation}


В (5.6), (5.7) $\Gm$ играет роль программного оператора, подобно


ОПП ${\Bbb A}$ (4.3). Между $\Gm$ и ${\Bbb A}$ существует тесная связь.





Через ${{\cal P}_0}({\bold D})$ обозначаем семейство всех множеств


$H \in {{\cal P}({\bold D})}$ таких, что


$\fo{t} \in [{{\tau}_0},{{\vt}_0}[$,


$\fo{{\omega}_1} \in \Om$, $\fo{u} \in {{\cal C}_0}({\omega}_1)$


\begin{equation}


((t,u) \in H) \Lo (\fo{{\omega}_2} \in {{\wt \Om}_t}({\omega}_1)


\ \fo{v} \in {{\cal C}_0}({\omega}_2) \cap {{\wt Z}_t}(u):


(t,v) \in H).


\end{equation}


Можно рассматривать (5.8) как ослабленный вариант наследственности


соответствующего множества; всегда ${{\cal P}_{\bold h}}({\bold D}) \subset


{{\cal P}_0}({\bold D})$ и


\begin{equation}


{\bigcap\limits_{H \in {\cal


H}}}{H} \in {{\cal P}_0}({\bold D})\ \;


\fo{\cal H} \in {{\cal P}'}({{\cal P}_0}({\bold D})).


\end{equation}


В дальнейшем полагаем выполненным





\begin{conds}


$\fo{t} \in [{{\tau}_0},{{\vt}_0}[$, $\fo{{\omega}_1} \in \Om$,


$\fo{{\omega}_2} \in {{\wt \Om}_t}({\omega}_1)$, $\fo{u} \in


{{\cal C}_0}({\omega}_1)$, $\fo{v} \in {{\cal C}_0}({\omega}_2) \cap


{{\wt Z}_t}(u)$,


$\fo{\omega} \in \Om:S((t,u),\omega) = S((t,v),\omega)$.


\end{conds}





\begin{notes}


Условие 5.1 является некоторым ослаблением условия 3.5.


Условия 3.1--3.5 и все условия настоящего раздела


весьма традиционны для систем управляемых дифференциальных уравнений


([28], [29] и др.); относительно ограничительными представляются


условия 3.1, 3.2. Для осуществления последних в нелинейных системах с


геометрическими ограничениями на выбор управлений, являющихся


более общими в сравнении с (3.2), следует использовать расширения с


применением скользящих режимов ([2], [3], [5], [28], [29] и др.).


\end{notes}





Возвращаясь к (5.8), отметим,


что  ${{\bold A}_\omega}(H) \in {{\cal P}_0}({\bold D})$


$\fo{H} \in {{\cal P}_0}({\bold D})$,


$\fo{\omega} \in \Om$.


С учетом (5.9) получаем ${{\Bbb A}(H)} \in {{\cal P}_0}({\bold D})$


$\fo{H} \in {{\cal P}_0}({\bold D})$.


Следовательно, ${{\Bbb A}^k}({\bold D}) \in {{\cal P}_0}({\bold D})$


$\fo{k} \in {{\cal N}_0}$.


Вновь используя (5.9), получаем $\up{\infty}{\Bbb A}({\bold D})


\in{{\cal P}_0}({\bold D})$.





Если ${{\omega}_1} \in \Om,\;{{\omega}_2} \in \Om$ и $t \in {I_0}$,


то ${({{\omega}_1}\Box{{\omega}_2})_t} \in {\Upsilon}^{I_0}$ определяем


в виде


\begin{equation}


(\fo{\xi} \in [{t_0},t]:{({{\omega}_1}\Box{{\omega}_2})_t}(\xi) \stac


{{\omega}_1}(\xi)) \& (\fo{\xi} \in\, ]t,{{\vt}_0}]:


{({{\omega}_1}\Box{{\omega}_2})_t}(\xi) \stac {{\omega}_2}(\xi)).


\end{equation}


В связи с (5.10)


постулируем в дальнейшем естественное


для теории управления [28], [36]


\pagebreak





\begin{conds}


$\fo{{\omega}_1} \in \Om$, $\fo{{\omega}_2} \in \Om$, $\fo{t} \in


[{{\tau}_0},{{\vt}_0}]:{({{\omega}_1}\Box{{\omega}_2})_t} \in \Om$.


\end{conds}





Итак, в дальнейшем склейки (5.10),


реализуемые при $t \in [{{\tau}_0},{{\vt}_0}]$,


не выводят за пределы $\Om$.


Пусть далее выполнены следующие два условия.





\begin{conds}


$\fo{{\omega}_1} \in \Om$, $\fo{{\omega}_2} \in \Om$, $\fo{t} \in


[{{\tau}_0},{{\vt}_0}[$, $\fo{h} \in


{{\cal C}_0}({({{\omega}_1}\Box{{\omega}_2})_t}):h \in


S((t,h),{{\omega}_2})$.


\end{conds}





\begin{conds}


$\fo{{\omega}_1} \in \Om$, $\fo{{\omega}_2} \in \Om$, $\fo{h} \in


{{\cal C}_0}({\omega}_1)$, $\fo{t} \in [{{\tau}_0},{{\vt}_0}[$


$$


S((t,h),{{\omega}_2}) \cap M \subset


S({z_0},{({{\omega}_1}\Box{{\omega}_2})_t}) \cap {{\wt Z}_t}(h).


$$


\end{conds}





\begin{notes}


Отметим, что в традиционном случае управляемых дифференциальных


уравнений типа рассматриваемых в [2], [4] справедливость последних


двух условий очевидна в свете известного полугруппового свойства решений.


В целом условия 5.1--5.4 имеют иной, в сравнении с условиями раздела


3, характер (см. в этой связи краткое изложение в [37], где введены


упомянутые группы условий). В частности, отметим, что в последующих


утверждениях настоящего раздела не предполагается выполнение важных


для теоремы 4.1 условий 3.2, 3.3 топологического характера. Следовательно,


двойственность итерационных методов, определяемая фактически уже в следующем


предложении, охватывает случаи, для которых непрямая итерационная


процедура раздела 4 может не обеспечивать построение требуемой


неподвижной точки ОПП.


\end{notes}





\begin{props}


Если $H \in {{\cal P}_0}({\bold D})$, то ${\Gm}({\Pi}(\cdot \mid


{z_0},H)) = {\Pi}(\cdot \mid {z_0},{{\Bbb A}(H)})$.


\end{props}





\begin{pf}


Пусть $H \in {{\cal P}_0}({\bold D})$,


$\alpha = {\Pi}(\cdot \mid {z_0},H)$, $\beta \stac {\Gm}(\alpha)$ и


$\delta \stac {\Pi}(\cdot \mid {z_0},{{\Bbb A}(H)})$. При


$\nu \in \Om$ имеем ${{\alpha}(\nu)}$ --- множество всех


$h \in {{\cal C}_0}(\nu)$ таких, что $\fo{t} \in [{{\tau}_0},{{\vt}_0}[\,:


(t,h) \in H$; в этих же условиях ${\delta}(\nu)$ --- множество всех


${\wt h} \in {{\cal C}_0}(\nu)$ таких, что


$(t,{\wt h}) \in {{\Bbb A}(H)}$ $\fo{t} \in [{{\tau}_0},{{\vt}_0}[$.


Фиксируем $\omega \in \Om$. Пусть $\vp \in


{{\beta}(\omega)}$. В частности, $\vp \in {{\alpha}(\omega)}$


(см. (5.3)). Следовательно, $\vp \in {{\cal C}_0}(\omega)$. Пусть


$\vp \notin {{\delta}(\omega)}$: для некоторого ${t_*} \in


[{{\tau}_0},{{\vt}_0}[$ имеем ${z_*} \stac ({t_*},\vp) \in H \setminus


{{\Bbb A}(H)}$. В силу (4.3) для некоторого ${{\omega}'} \in


\Om$ выполняется ${\Pi}({\omega}' \mid {z_*},H) = \emp$,


\begin{equation}


{{\omega}''} \stac (\omega\Box{{\omega}'})_{t_*} \in


{{\wt \Om}_{t_*}}(\omega).


\end{equation}


С учетом (5.3), (5.11) имеем ${{\wt Z}_{t_*}}(\vp) \cap


{{\alpha}({\omega}'')} \ne \emp$. Пусть $\psi \in


{{\wt Z}_{t_*}}(\vp) \cap {{\alpha}({\omega}'')}$. В частности,


$\psi \in {{\cal C}_0}({\omega}'')$. В силу условия 5.3 и (5.11)


$\psi \in S(({t_*},\psi),{\omega}')$.


Из условия 5.1 следует $S({z_*},{\omega}') = S(({t_*},\psi),{\omega}')$.


В итоге $\psi \in S({z_*},{\omega}') \cap M$. По свойствам $\alpha$


имеем $(t,\psi) \in H$ $\fo{t} \in [{{\tau}_0},{{\vt}_0}[$.


Как следствие, $\psi \in {\Pi}({\omega}' \mid {z_*},H)$, что противоречит


выбору ${{\omega}'}$. Итак,


$\vp \in {{\delta}(\omega)}$. Вложение ${{\beta}(\omega)} \subset


{{\delta}(\omega)}$ установлено.


Пусть $u \in {{\delta}(\omega)}$. Тогда $u \in {{\cal C}_0}(\omega)$ и


\begin{equation}


(t,u) \in {{\Bbb A}(H)}\ \; \fo{t} \in [{{\tau}_0},{{\vt}_0}[.


\end{equation}


Из (5.12) следует, что $(t,u) \in H$  $\fo{t} \in [{{\tau}_0},{{\vt}_0}[$.


В итоге $u \in {{\alpha}(\omega)}$. Пусть


${t^*} \in [{{\tau}_0},{{\vt}_0}[$ и ${{\omega}^*} \in


{{\wt \Om}_{t^*}}(\omega)$. Рассмотрим ${z^*} \stac ({t^*},u) \in


{{\Bbb A}(H)}$. Тогда (см. (4.3)) ${\Pi}({{\omega}^*} \mid


{z^*},H) \ne \emp$. Пусть $v \in {\Pi}({{\omega}^*} \mid {z^*},H)$.


Тогда $v \in S({z^*},{{\omega}^*}) \cap M$ и


\begin{equation}


\fo{t} \in [{t^*},{{\vt}_0}[\,:(t,v) \in H.


\end{equation}


Заметим здесь же, что ${({\omega}\Box{{\omega}^*})}_{t^*}={{\omega}^*} \in


\Om$ и в силу условия 5.4


\begin{equation}


S({z^*},{{\omega}^*}) \cap M \subset S({z_0},{{\omega}^*}) \cap


{{\wt Z}_{t^*}}(u).


\end{equation}


Из (5.13), (5.14) получаем $v \in S({z_0},{{\omega}^*}) \cap


{{\wt Z}_{t^*}}(u)$. Тогда, в частности,


\begin{equation}


v \in {{\cal C}_0}({\omega}^*) \cap {{\wt Z}_{t^*}}(u).


\end{equation}


Пусть $\xi \in [{{\tau}_0},{t^*}[$.


Поскольку $\xi \in [{{\tau}_0},{{\vt}_0}[$, то (см. (5.8))


\begin{equation}


((\xi,u) \in H) \Lo (\fo{\wt \omega} \in {{\wt \Om}_\xi}(\omega),


\ \fo{q} \in {{\cal C}_0}({\wt \omega}) \cap {{\wt Z}_\xi}(u):


(\xi,q) \in H).


\end{equation}


Но $[{t_0},\xi] \subset [{t_0},{t^*}]$, ${{\omega}^*} \in


{{\wt \Om}_\xi}(\omega)$ и в силу (5.15) $v \in


{{\cal C}_0}({{\omega}^*}) \cap {{\wt Z}_\xi}(u)$.


Из (5.16) имеем $(\xi,v) \in H$.


Установлено, что $(t,v) \in H$ $\fo{t} \in [{{\tau}_0},{t^*}[$.


С учетом (5.13), (5.15) $v \in {{\alpha}({\omega}^*)} \cap


{{\wt Z}_{t^*}}(u)$. Следовательно,


\begin{equation}


\fo{t} \in [{{\tau}_0},{{\vt}_0}[,\ \fo{\wt \omega} \in


{{\wt \Om}_t}(\omega):{{\alpha}({\wt \omega})} \cap


{{\wt Z}_t}(u) \ne \emp;


\end{equation}


так как ${{\wt \Om}_{{\vt}_0}}(\omega) = \{\omega\}$, из


(5.3) и (5.17) имеем $u \in {{\beta}(\omega)}$. Итак,


${{\delta}(\omega)} \subset {{\beta}(\omega)}$.


\end{pf}





Индуктивно доказывается


${{\Gm}^k}({\cal C}_0) = {\Pi}(\cdot \mid


{z_0},{{\Bbb A}^k}({\bold D}))$\;


$\fo{k} \in {{\cal N}_0}$.





\begin{thms}


Справедливо равенство ${{\Gm}^\infty}({\cal C}_0) =


{\Pi}(\cdot \mid {z_0},\up{\infty}{\Bbb A}({\bold D}))$.


\end{thms}





\begin{pf}


Пусть $\omega \in \Om$. Из (5.6), предложения 5.1


и теоремы 4.1 имеем


\begin{equation}


{\Pi}(\omega \mid {z_0},\up{\infty}{\Bbb A}({\bold D})) \subset


{{\Gm}^\infty}({\cal C}_0)(\omega).


\end{equation}


Пусть $\vp \in {{\Gm}^\infty}({\cal C}_0)(\omega)$. В силу предложения 5.1


$\fo{k} \in {{\cal N}_0}:\vp \in {\Pi}(\omega \mid


{z_0},{{\Bbb A}^k}({\bold D}))$.


Как следствие, $\vp \in {{\cal C}_0}(\omega)$ таково, что


$(t,\vp) \in {{\Bbb A}^k}({\bold D})$\;


$\fo{k} \in {{\cal N}_0}$, $\fo{t} \in [{{\tau}_0},{{\vt}_0}[$.


Поэтому при $t \in [{{\tau}_0},{{\vt}_0}[$ имеем


$(t,\vp) \in \up{\infty}{\Bbb A}({\bold


D})$. Тогда $\vp \in {\Pi}(\omega \mid


{z_0},\up{\infty}{\Bbb A}({\bold D}))$, чем и завершается


доказательство вложения, противоположного (5.18).


\end{pf}





\section{Многозначные квазистратегии и метод итераций}





Вернемся к


МО (5.5), являющемуся предельно широкой неупреждающей


реакцией на развитие непредсказуемых факторов.


Из [22], [23], [31] следует


$(\na)[{\cal C}_0] \sqsubseteq {{\Gm}^\infty}({\cal C}_0)$.


Постулируем





\begin{conds}


$\fo{\omega} \in \Om:S({z_0},\omega) \in {\frak K}$,


\end{conds}





\noindent{}из которого следует


\begin{equation}


\fo{\omega} \in \Om:{{\cal C}_0}(\omega) \in {\frak K}.


\end{equation}


Далее полагаем, что ${\bold t}$


--- топология множества $X$, причем


\begin{equation}


(X,{\bold t})


\end{equation}


есть хаусдорфово ТП,


и считаем выполненным





\begin{conds}


Топология ${\cal T}$ индуцирована на ${\Bbb C}$ ([25], c.\,77)


тихоновским произведением


\begin{equation}


({X^{I_0}},{{\otimes}^{I_0}}({\bold t}))


\end{equation}


экземпляров ТП (6.2) с индексным множеством $I_0$, где


через ${{\otimes}^{I_0}}({\bold t})$ обозначена


топология множества $X^{I_0}$.


\end{conds}





Так как ${\cal T}$ --- топология поточечной сходимости


на ${\Bbb C}$ ([25], c.\,283), соответствующая оснащению (6.2) множества


$X$, а $({\Bbb C},{\cal T})$ есть п/п ТП (6.3), то


([22], [23], [31]) из (6.1) имеем


\begin{equation}


{{\Gm}^\infty}({\cal C}_0) = (\na)[{\cal C}_0] \in {{\frak K}^\Om}.


\end{equation}


Из теоремы 5.1 и (6.4) получаем, в частности, равенство


\begin{equation}


(\na)[{\cal C}_0] = {\Pi}(\cdot \mid {z_0},\up{\infty}{\Bbb A}({\bold


D})).


\end{equation}


Если значения МО (5.5) непусты, то в


(6.4), (6.5) определяется экстремальная по порядку


квазистратегия, реализующая движения из


$M$. Пусть


${(\DOM)[\alpha]}


\stac \{\omega \in \Om \mid {{\alpha}(\omega)} \ne \emp\}$


$\fo{\alpha} \in {\Bbb M}(\Om,{\Bbb C})$.





\begin{props}


Если ${(\DOM)[(\na)[{\cal C}_0]]} = \Om$, то ${z_0} \in


\up{\infty}{\Bbb A}({\bold D})$.


\end{props}





\begin{pf}


Пусть $(\DOM)[(\na)[{\cal C}_0]] = \Om$, но


${z_0} \notin \up{\infty}{\Bbb A}({\bold D})$;


тогда ${\Bbb K} \stac \{k \in {{\cal N}_0}


\mid {z_0} \notin {{\Bbb A}^k}({\bold D})\} \in {{\cal P}'({\cal N})}$,


$n \stac {{\inf}({\Bbb K})} \in {\Bbb K}$,


$n - 1 \in {{\cal N}_0} \setminus {\Bbb K}$ и


${z_0} \in {{\Bbb A}^{n-1}}({\bold D}) \setminus {{\Bbb A}^n}({\bold D})$.


Из (4.3) имеем для некоторого $\omega \in \Om$:


${\Pi}(\omega \mid {z_0},{{\Bbb A}^{n-1}}({\bold D})) = \emp$; тогда


${\Pi}(\omega \mid {z_0},


\up{\infty}{\Bbb A}({\bold D})) = \emp$. В силу (6.5)


$(\na)[{\cal C}_0](\omega) = \emp$.


\end{pf}





\begin{props}


Пусть выполнены условия $3.2$, $3.3$ и ${z_0} \in


\up{\infty}{\Bbb A}({\bold D})$. Тогда


$$


(\DOM)[(\na)[{\cal C}_0]] = \Om.


$$


\end{props}





\begin{pf}


Поскольку (см. теорему 4.1) ${z_0} \in


{\Bbb A}(\up{\infty}{\Bbb A}({\bold D}))$, то в силу (4.3)


${\Pi}(\omega \mid {z_0},


\up{\infty}{\Bbb A}({\bold D})){\ne}\emp$ $\fo{\omega} \in \Om$. Теперь


учитываем (6.5).


\end{pf}





\begin{thebibliography}{99}





\bibitem{1} Красовский Н.Н. {\em Игровые задачи о встрече движений}. --


М.: Наука, 1970. -- 420 c.





\bibitem{2} Красовский Н.Н., Субботин А.И. {\em Позиционные дифференциальные


игры}. -- М.: Наука, 1974. -- 456 с.





\bibitem{3} Красовский Н.Н. {\em Управление динамической системой.


Задача о минимуме гарантированного результата}. -- М.: Наука,


1985. -- 518 с.





\bibitem{4} Субботин А.И. {\em Минимаксные неравенства и уравнения


Гамильтона--Якоби}. -- М.: Наука, 1991. -- 215 с.





\bibitem{5} Субботин А.И., Ченцов А.Г. {\em Оптимизация гарантии в


задачах управления}. -- М.: Наука, 1981. -- 287 с.





\bibitem{6} Айзекс Р. {\em Дифференциальные игры}. -- М.: Мир, 1967.


-- 479 с.





\bibitem{7} Ченцов А.Г. {\em О структуре одной игровой задачи сближения}


// ДАН СССР. -- 1975. -- Т.\,224. -- \No\,6. -- C.\,1272--1275.





\bibitem{8} Ченцов А.Г. {\em К игровой задаче наведения} //


ДАН СССР. -- 1976. -- Т.\,226. -- \No\,1. -- C.\,73--76.





\bibitem{9} Ченцов А.Г. {\em К игровой задаче наведения с информационной


памятью} // ДАН СССР. -- 1976. -- Т.\,227. -- \No\,2. -- C.\,306--308.





\bibitem{10} Ченцов А.Г. {\em Об игровой задаче сближения в заданный момент


времени} // Матем. сб. -- 1976. -- Т.\,99. -- \No\,3. -- C.\,394--420.





\bibitem{11} Чистяков С.В. {\em К решению игровых задач преследования} //


ПMM. -- 1977. -- Т.\,41. -- \No\,5. -- C.\,825--832.





\bibitem{12} Ченцов А.Г. {\em Метод программных итераций для дифференциальной


игры сближения-укло\-не\-ния}. -- Ин-т матем. и мех.


УНЦ АН СССР, 1979. -- 102 c. -- Деп. в ВИНИТИ, \No\,1933-79. 





\bibitem{13} Ченцов А.Г. {\em Об игровой задаче наведения к заданному


моменту времени} // Изв. АН СССР. Сер. матем. --


1978. -- Т.\,42. -- \No\,2. -- С.\,455--467.





\bibitem{14} Ченцов А.Г. {\em Итерационная программная конструкция для


дифференциальной игры с фиксированным моментом окончания} //


ДАН СССР. -- 1978. -- T.\,240. -- \No\,1. -- C.\,36--39.





\bibitem{15} Дятлов В.П., Ченцов А.Г. {\em Монотонные итерации множеств и их


приложения к игровым задачам управления} // Кибернетика. -- 1987. --


\No\,2. -- C.\,92--99.





\bibitem{16} Saint-Pierre P. {\em Approximation of the viability


kernel} // Appl. math. optim. -- 1992. -- V.\,29. -- P.\,187--209.





\bibitem{17} Cardaliaguet P., Quincampoix M., Saint-Piarre P.


{\em Some algorithms for differential games with two players and one


target} // Cahiers de math. de la decision. -- 1993. --


\No\,9343. -- 19~p.





\bibitem{18} Субботин А.И., Ченцов А.Г. {\em Итерационная процедура для


построения минимаксных и вязкостных решений уравнений


Гамильтона--Якоби} // Докл. PАН. -- 1996. -- Т.\,348. --


\No\,6. -- C.\,736--739.





\bibitem{19} Субботин А.И., Ченцов А.Г. {\em Итерационная процедура


построения минимаксных и вязкостных решений уравнений


Гамильтона--Якоби и ее обобщения} // Тр. МИАН. -- 1999. -- Т.\,224. --


C.\,311--334.





\bibitem{20} Ченцов А.Г. {\em Метод программных итераций в классе


конечно-аддитивных управлений-мер} // Дифференц. уравнения. --


1997. -- \No\,11. -- C.\,1528--1536.





\bibitem{21} Ченцов А.Г. {\em Итерационная реализация неупреждающих


многозначных отображений} // Докл. PАН. -- 1997. --


Т.\,357. -- \No\,5. -- C.\,595--598.





\bibitem{22} Ченцов А.Г. {\em К вопросу о параллельной версии


абстрактного аналога метода программных итераций} // Докл.


PАН. -- 1998. -- Т.\,362. -- \No\,5. -- С.\,602--605.





\bibitem{23} Ченцов А.Г. {\em К вопросу об итерационной реализации


неупреждающих многозначных отображений} // Изв. вузов.


Математика. -- 2000. -- \No\,3. -- C.\,66--76.





\bibitem{24} Эдвардс Р. {\em Функциональный анализ. Теория


и приложения}. -- М.: Мир, 1969. -- 1071 с.





\bibitem{25} Келли Дж.Л. {\em Общая


топология}. -- М.: Наука, 1981. -- 433 c.





\bibitem{26} Александрян Р.А., Мирзаханян Э.А. {\em Общая топология}.


Учеб. пособие. -- М.: Высш. школа, 1979. -- 336~с.





\bibitem{27} Энгелькинг Р. {\em Общая топология}. --


М.: Мир, 1986. -- 751 с.





\bibitem{28} Варга Дж. {\em Оптимальное управление дифференциальными и


функциональными уравнениями}. -- М.: Наука, 1977. -- 624 с.





\bibitem{29} Гамкрелидзе Р.В. {\em Основы оптимального управления}. --


Тбилиси: Изд-во Тбилисск. ун-та, 1975. -- 229 с.





\bibitem{30} Chentsov A.G. {\em Finitely additive measures and


relaxations of extremal problems}.


-- New York--London--Moscow:


Plenum Publishing Corporation, 1996. -- 244 p.





\bibitem{31} Chentsov A.G. {\em Nonanticipating selectors of set-valued


mappings and iterated procedures} // Funct. Diff.


Equat. -- 1999. -- V.6. -- \No\,3--4. -- P.\,249--274.





\bibitem{32} Elliott R.J., Kalton N.J. {\em The existence of value in


differential games} // Memoirs of the Amer. Math. Soc. -- 1972.


-- \No\,126. -- 67 p.





\bibitem{33} Krasovskii N.N., Chentsov A.G. {\em On the design of differential


games}, I // Probl. Control and Inform. Theory. -- 1977. -- V.\,6. --


\No\,5--6. -- P.\,381--395.





\bibitem{34} Krasovskii N.N., Chentsov A.G. {\em On the design of differential


games}, II // Probl. Control and Inform. Theory. -- 1979. -- V.\,8.


-- \No\,1. -- P.\,3--11.





\bibitem{35} Roxin E. {\em Axiomatic approach in differential games} //


J. Optimiz. Theory and Appl. -- 1969. -- V.\,3. -- \No\,3. -- P.\,153--163.





\bibitem{36} Понтрягин Л.С., Болтянский В.Г., Гамкрелидзе Р.В.,


Мищенко Е.Ф. {\em Математическая теория оптимальных процессов}.


-- М.: Наука, 1976. -- 392 с.





\bibitem{37} Ченцов А.Г. {\em Топологический вариант метода программных


итераций в абстрактной задаче управления}


// Докл. PАН. -- 2001. -- Т.\,376. -- \No\,3. --


C.\,311-314.








\end{thebibliography}





\vskip 2cm





\post{\it Институт математики и механики}{\it Поступила}


\post{\it Уральского отделения}{13.12.2000}


\post{\it Российской Академии наук}{}





\label{end}


\end{document}


